\chapter{Fonctions serveur (Edge Functions IA)}

Toute la logique faisant appel à l'\textbf{intelligence artificielle} est
déportée dans des \textbf{Edge Functions Deno} hébergées par Supabase
(\filepath{supabase/functions/}). Ce choix garantit que les clés d'API ne sont
\textbf{jamais} exposées au navigateur~: elles sont lues côté serveur via
\texttt{Deno.env.get(...)}.

\section{Fournisseurs et modèles d'IA}
\begin{longtable}{p{5cm}p{9.5cm}}
\toprule \textbf{Fournisseur} & \textbf{Usage} \\ \midrule \endhead
\textbf{Lovable AI Gateway} & Passerelle principale vers Google Gemini
(\texttt{google/gemini-2.5-flash}, \texttt{google/gemini-3-flash-preview}).
Clé \texttt{LOVABLE\_API\_KEY}. \\
\textbf{Groq} & Modèle de repli rapide \texttt{llama-3.3-70b-versatile}.
Clé \texttt{GROQ\_API\_KEY}. \\
\textbf{Google Gemini (direct)} & Génération en masse via \texttt{GEMINI\_API\_KEY} /
\texttt{GEMINI\_API\_KEY\_2}. \\
\textbf{OpenRouter / Cloudflare AI} & Clés additionnelles disponibles
(\texttt{OPENROUTER\_API\_KEY}, \texttt{CLOUDFLARE\_AI\_API\_KEY}). \\
\bottomrule
\end{longtable}

\begin{notebox}[\faLock\ Stratégie de repli]
Plusieurs fonctions tentent d'abord Gemini (via la passerelle) puis basculent
automatiquement sur Groq (Llama) en cas d'erreur, afin de garantir la
disponibilité du service.
\end{notebox}

\section{Catalogue des fonctions}
\begin{longtable}{p{5cm}p{9.5cm}}
\toprule \textbf{Fonction} & \textbf{Rôle} \\ \midrule \endhead
\texttt{lovable-chat} & Assistant pédagogique IA (chatbot). Réponses
contextualisées et strictement encadrées par le programme de mathématiques.
Gemini avec repli Llama. \\
\texttt{generate-adaptive-content} & Génère quiz / exercices / révisions
calibrés au niveau composite de l'élève (Gemini 3 flash). \\
\texttt{generate-lesson-comment} & Produit un commentaire IA personnalisé en
arabe après une session d'étude (forces/lacunes, conseils). \\
\texttt{generate-remediation} & Génère des activités de remédiation ciblées
\emph{uniquement} sur les lacunes détectées de l'élève. \\
\texttt{generate-chapter-quizzes} & Génère des quiz pour tout un chapitre
(Gemini + repli Llama). \\
\texttt{generate-chapter-revision} & Génère des fiches de révision de chapitre. \\
\texttt{generate-parent-report} & Construit le rapport de progression destiné
aux parents (synthèse + recommandations IA). \\
\texttt{scheduled-parent-reports} & Tâche planifiée (pg\_cron) générant
quotidiennement les rapports pour chaque liaison parent--enfant active. \\
\texttt{enrich-chapter-lessons} & Enrichit le contenu des leçons d'un chapitre. \\
\texttt{bulk-generate-lesson-content} & Génère en masse 10 exercices + 10 quiz
pour une leçon et les insère en base. \\
\texttt{bulk-generate-all-terminale} & Génération par lot pour toutes les leçons
de Terminale Sciences. \\
\texttt{bulk-gen-terminale-gemini} & Variante de génération en masse via l'API
Gemini directe (\texttt{GEMINI\_API\_KEY\_2}). \\
\texttt{delete-user-account} & Suppression complète et sécurisée d'un compte
utilisateur (privilèges service\_role). \\
\bottomrule
\end{longtable}

\section{Secrets configurés}
Les secrets suivants sont disponibles côté Edge Functions (jamais exposés au
client)~:
\begin{itemize}
  \item IA~: \texttt{LOVABLE\_API\_KEY}, \texttt{GROQ\_API\_KEY},
        \texttt{GEMINI\_API\_KEY}, \texttt{GEMINI\_API\_KEY\_2},
        \texttt{OPENROUTER\_API\_KEY}, \texttt{CLOUDFLARE\_AI\_API\_KEY},
        \texttt{CLOUDFLARE\_ACCOUNT\_ID}.
  \item Supabase~: \texttt{SUPABASE\_URL}, \texttt{SUPABASE\_DB\_URL},
        \texttt{SUPABASE\_SERVICE\_ROLE\_KEY}, \texttt{SUPABASE\_ANON\_KEY},
        \texttt{SUPABASE\_PUBLISHABLE\_KEY}, \texttt{SUPABASE\_JWKS}\dots
\end{itemize}
